\section{Difference in means problem: A counterfactual exercise} 
\label{sec:counterfactual_example}

Before introducing the OBD, we first illustrate the index-number problem by framing it as a counterfactual exercise as in \citet{fortinLemieuxFirpo2011}. 
We consider two groups, $H$ and $K$, and an outcome $Y$ we wish to compare across groups, together with observed characteristics $X$ (i.e., covariates that are relevant predictors of $Y$). 
Although many counterfactuals can be constructed, a simple and intuitive pair arises by using the other group's outcome--covariate relationship as a bridge. In particular, one could ask one of the following two questions:

\begin{center}
\emph{“What would Group $K$’s mean outcome ($Y$) be if its observed characteristics ($X$) were evaluated under Group $H$’s outcome–covariate relationship ($Y|X$)?”}

\emph{“What would Group $H$’s mean outcome ($Y$) be if its observed characteristics ($X$) were evaluated under Group $K$’s outcome–covariate relationship ($Y|X$)?”}
\end{center}

Although both questions are sensible, they may not lead to the same answer. 
Which one to ask is rarely obvious and depends on the goal of the analysis.
For instance, drawing from the classical gender wage-gap example \citep{oaxaca1973, blinder1973}, if $Y$ denotes wages and $X$ education, experience, occupation, etc., these questions ask \emph{What would women’s ($K$) mean wage be if their observed characteristics were remunerated at men’s ($H$) returns?} And \emph{What would men’s ($H$) mean wage be if their observed characteristics were remunerated at women’s ($K$) returns?}

In a similar spirit, from our motivating health example, if $Y$ is systolic blood pressure (SBP) and $X$ is body mass index (BMI), this question asks: \emph{What would City $K$’s mean SBP be if its observed BMI were evaluated under City $H$’s SBP–BMI relationship?} And \emph{What would City $H$’s mean SBP be if its observed BMI were evaluated under City $K$’s SBP–BMI relationship?} 

Which counterfactual is more appropriate depends on the question of interest and carries distinct interpretations and policy implications.

All the above questions define at least two counterfactual means \(\E\!\big[Y_{K}^{C(H)}\big],\) and \(\E\!\big[Y_{H}^{C(K)}\big],\) where \(Y_{K}^{C(H)}\) denotes the counterfactual outcome for Group $K$ if its observed characteristics were mapped through Group $H$’s conditional outcome rule and $Y_{H}^{C(K)}$ denotes the counterfactual outcome for Group $H$ if its observed characteristics $X_H$ were mapped through Group $K$’s conditional outcome rule. Note that we use the superscript \(C(\cdot)\) to indicate counterfactual outcomes. 

Using these counterfactuals, the difference in mean outcomes between groups can be expressed as the distance between each observed mean and its corresponding counterfactual:
\begin{align}
    \label{eq:OB_H_Counterfactual}
    \E_H[Y] &- \E_K[Y] = \\ 
    &\underbrace{\Big(\E_H[Y] - \E[Y_{K}^{C(H)}]\Big)}_{\text{Explained}:= E^{(H)}} \nonumber 
     + \underbrace{\Big(\E[Y_{K}^{C(H)}] - \E_K[Y]\Big)}_{\text{Unexplained} := U^{(H)} }.
\end{align} 
Alternatively, 
\begin{align}
    \label{eq:OB_K_Counterfactual}
    \E_H[Y] & - \E_K[Y]= \\
    & \underbrace{\Big(\E_H[Y] - \E[Y_{H}^{C(K)}]\Big)}_{\text{Unexplained}:= U^{(K)}} \nonumber 
     + \underbrace{\Big(\E[Y_{H}^{C(K)}] - \E_K[Y]\Big)}_{\text{Explained}:= E^{(K)}}.
\end{align}

Certainly, the choice of the counterfactual is not unique or canonical, and the two decompositions in \cref{eq:OB_H_Counterfactual,eq:OB_K_Counterfactual} are mathematically equivalent, as they both sum to the same total difference in means.
OBD is nothing more than a specific implementation of these two counterfactual decompositions under the assumption of a linear model for the outcome-generating process.
